\documentclass[11pt]{article}
\input{packages.tex}
\usepackage{ulem}
\DeclareMathOperator{\sign}{sign}
\DeclareMathOperator{\Imm}{Im}
\DeclareMathOperator{\Ker}{Ker}
\DeclareMathOperator{\rg}{rg}

% Настройка колонтитулов
\pagestyle{fancy}
\fancyhf{} % Очистить все поля
\fancyhead[L]{Всё приходит к тому, кто умеет ждать \sout{Подлипского в КПМ}}
\fancyhead[R]{\href{https://t.me/ceagest}{Автор}}
\fancyfoot[C]{\thepage}               % Справа внизу

\setlength{\headheight}{13.6pt}

\hypersetup{
    colorlinks=true,
    linkcolor=blue,
    citecolor=blue,
    filecolor=blue,
    urlcolor=blue,
}

\begin{document}

\begin{titlepage}
    \centering
    \vspace*{\fill}  % заполняет пространство сверху

    {\Huge Лекция 6. Линейные преобразования. Инвариантные подпространства. Собственные значения и подпространства. Собственные векторы. Характеристические уравнения
    \par}

    \vspace*{\fill}  % заполняет пространство снизу
\end{titlepage}

\setcounter{page}{1}

\section{Дополнение к предыдущей лекции}

\begin{definition}[Обратное отображение]
    Пусть дано линейное отображение $A \colon \mathcal{L} \to \widetilde{\mathcal{L}}$. Линейное отображение $B \colon \widetilde{\mathcal{L}} \to \mathcal{L}$ назовём \textbf{обратным} для $A$ и будем обозначать 
    $A^{-1}$, если $BA = E$ и $AB = \widetilde{E}$, где $E$ и $\widetilde{E}$ --- тождественные преобразования пространств $\mathcal{L}$ и $\widetilde{\mathcal{L}}$ соответственно. Иначе говоря, для любых $x \in \mathcal{L}$ и $y \in \widetilde{\mathcal{L}}$ должно быть выполнено:
    \[B(A(x)) = x, \quad A(B(y)) = y\]
\end{definition}

\begin{theorem}
    Линейное отображение $A \colon \mathcal{L} \to \widetilde{\mathcal{L}}$ имеет обратное 
    тогда и только тогда, когда оно является изоморфизмом.
\end{theorem}

\begin{proof}
    Пусть $\{e_1, \ldots, e_n\}$ и $\{f_1, \ldots, f_m\}$ --- базисы в $\mathcal{L}$ и $\widetilde{\mathcal{L}}$ соответственно. Пусть $A$ --- матрица отображения $A$ в этих базисах.
    \\
    $(\Longleftarrow):$ Пусть $A$ --- изоморфизм. Тогда $A$ --- невырожденная квадратная матрица, а значит она имеет обратную матрицу $A^{-1}$. Рассмотрим отображение $B \colon \widetilde{\mathcal{L}} \to \mathcal{L}$, 
    определяемое матрицей $A^{-1}$ в базисах $f$ и $e$. Оно и является искомым обратным отображением к $A$, так как:
    \[\forall x \in \mathcal{L} \; \; \forall y \in \widetilde{\mathcal{L}} \; \; \hookrightarrow \; \; B(A(x)) = (BA)x = (A^{-1}A)x = x, \quad A(B(y)) = (AB)y = (A A^{-1})y = y\]
    $(\Longrightarrow):$ Пусть $A$ не является изоморфизмом. Тогда из доказанного в предыдущей лекции вытекает, что либо $\rg A < m$, либо $\rg A < n$. В первом случае отображение $A$ не будет являться сюръекцией, а значит найдётся вектор $u$ из $\widetilde{\mathcal{L}}$, не 
    принадлежащий $\Imm A$. Если существует обратное отображение, мы приходим к противоречию:
    \[u = A(A^{-1}(u)) \in \Imm A\]
    Во втором случае отображение $A$ не будет являться инъекцией, откуда по доказанному в предыдущей лекции 
    размерность $\Ker A$ не равна нулю, что эквивалентно существованию ненулевого вектора $z \in \mathcal{L}$, лежащего в ядре $A$. Если существует $A^{-1}$, то вновь получаем противоречие:
    \[z = A^{-1}(A(z)) = A^{-1}(0) = 0\]
\end{proof}

\begin{remark}
    Из доказательства выше видно, что матрицей обратного отображения в базисах $f$ и $e$ является $A^{-1}$.
\end{remark}

\section{Линейные преобразования}

\subsection{Воспоминания}

Сформулируем ещё раз определение матрицы линейного отображения для частного его случая --- линейного преобразования.

\begin{definition}
    Матрицей линейного преобразования $A \colon \mathcal{L} \to \mathcal{L}$ в базисе $\{e_1, \ldots, e_n\}$ 
    называется матрица, столбцы которой --- координатные столбцы векторов $A(e_1), \ldots, A(e_n)$ в базисе $e$.
\end{definition} \leavevmode
\\
Если $S$ --- матрица перехода между базисами $e$ и $e'$ в пространстве $\mathcal{L}$, то по доказанному в предыдущей лекции (да да, мы будем очень часто ссылаться на замечательную пятую лекцию) матрица преобразования $A$ в базисе $e'$ есть 
\[A' = S^{-1} A S\]

\subsection{Умножение преобразований}

\begin{definition}
    Пусть $A$ и $B$ --- линейные преобразования пространства $\mathcal{L}$. Будем говорить, что $A$ и $B$ \textbf{перестановочны} или \textbf{коммутируют}, если $AB = BA$.
\end{definition}

\begin{definition}
    Пусть $A$ --- линейное преобразование пространства $\mathcal{L}$. Определим целую положительную степень $A$ по индукции соотношением $A^{k} = A A^{k - 1}$, положив базой 
    индукции $A^{1} = A$ и $AA = A^{2}$. Нулевой степенью преобразования $A$ по определению считают тождественное преобразование $E$ пространства $\mathcal{L}$, то есть $A^{0} := E$.
\end{definition}

\begin{definition}
    Линейное преобразование $B$, представленное как линейная комбинация целых неотрицательных степеней преобразования $A$
    \[\alpha_0 E + \alpha_1 A + \ldots + \alpha_n A^n,\]
    называется \textbf{многочленом от преобразования} $A$ или, точнее, значением многочлена 
    \[p(t) = \alpha_0 + \alpha_1 t + \ldots + \alpha_n t^n\]
    на преобразовании $A$, и обозначается $p(A)$.
\end{definition}

\begin{remark}
    Легко убедиться, что любой многочлен от $A$ перестановочен с $A$ и что любые два многочлена от $A$ перестановочны. Эти факты нам пригодятся для секции с собственными пространствами.
\end{remark}

\subsection{Инвариантные подпространства}

\begin{definition}
    Пусть $\mathcal{L}_1$ --- подпространство линейного пространства $\mathcal{L}$. Будем говорить, что подпространство $\mathcal{L}_1$ \textbf{инвариантно} относительно линейного преобразования $A$, если 
    \[\forall x \in \mathcal{L}_1 \hookrightarrow A(x) \in \mathcal{L}_1 \iff A(\mathcal{L}_1) \subset \mathcal{L}_1\]
\end{definition}

\begin{remark}
    Всё пространство $\mathcal{L}$, рассматриваемое как подпространство, и нулевое подпространство инвариантны относительно любого линейного преобразования.
\end{remark}

\begin{remark}
    Каждое подпространство является инвариантным относительно тождественного и нулевого преобразований.
\end{remark}

\begin{lemma}
    Пусть $A$ --- линейное преобразование пространства $\mathcal{L}$. Тогда подпространства $\Ker A$ и $\Imm A$ являются инвариантными относительно $A$.
\end{lemma}

\begin{proof}
    Подпространство $\Ker A$ инвариантно относительно $A$, так как:
    \[\forall x \in \Ker A \hookrightarrow A(x) = 0 \in \Ker A\]
    С подпространством $\Imm A$ всё совсем тривиально:
    \[\forall x \in \Imm A \hookrightarrow A(x) \in \Imm A\]
\end{proof}

\begin{theorem}
    Матрица $A$ линейного преобразования $A$ пространства $\mathcal{L}$ является \textbf{клеточно-треугольной}, то есть имеет вид
    \begin{equation*}
        A = 
        \left(\begin{array}{ c | c }
            A_1 & A_2 \\
            \hline
            O & A_3
        \end{array}\right)
        ,
    \end{equation*}
    тогда и только тогда, когда подпространство $\mathcal{L}_1 = \left\langle e_1, \ldots, e_k \right\rangle$, где $\{e_1, \ldots, e_n\}$ --- базис в $\mathcal{L}$, инвариантно относительно преобразования $A$. Здесь $A_1$, $A_2$ и $A_3$ --- некоторые матрицы размеров $k \times k$, $k \times (n - k)$ и $(n - k) \times (n - k)$ соответственно.
\end{theorem}

\begin{proof}
    $(\Longleftarrow):$ Пусть подпространство $\mathcal{L}_1 = \left\langle e_1, \ldots, e_k \right\rangle$ инвариантно относительно преобразования $A$. Выберем произвольно $x \in \mathcal{L}_1$. Так как $x$ лежит в подпространстве с базисом $\{e_1, \ldots, e_k\}$, то все его компоненты при векторах $\{e_{k + 1}, \ldots, e_n\}$ равны нулю:
    \[x = x_1 e_1 + x_2 e_2 + \ldots + x_k e_k + 0 \cdot e_{k + 1} + \ldots + 0 \cdot e_n\]
    При этом $\mathcal{L}_1$ инвариантно относительно $A$, а значит $A(x) \in \mathcal{L}_1$, откуда из соображений, аналогичных вышеописанным, получаем:
    \[A(x) = x_1' e_1 + x_2' e_2 + \ldots + x_k' e_k + 0 \cdot e_{k + 1} + \ldots + 0 \cdot e_n\]
    Подставляя теперь вместо $x$ поочерёдно базисные векторы $e_1, \ldots, e_k$, получим по определению матрицы линейного преобразования следующий вид матрицы $A$ (здесь для удобства за $a_{ij}$ обозначен коэффициент $x_i'$ в разложении $A(e_j)$ после подстановки $e_j$):
    \begin{equation*} \label{eq:th:2.1} \tag{$\ast$}
        A = 
        \left(\begin{array}{ c | c }
            \begin{matrix}
                a_{11} & a_{12} & \ldots & a_{1k} \\
                a_{21} & a_{22} & \ldots & a_{2k} \\
                \vdots & \vdots & \ddots & \vdots \\
                a_{k1} & a_{k2} & \ldots & a_{kk}
            \end{matrix} & A_2 \\
            \hline
            O & A_3
        \end{array}\right)
    \end{equation*}
    Легко видеть, что это и есть требуемый вид матрицы.
    \\
    $(\Longrightarrow):$ Пусть матрица $A$ линейного преобразования $A$ имеет тот самый вид (\ref{eq:th:2.1}). Тогда вновь из её определения вытекает:
    \begin{equation*}
        \begin{cases}
            A(e_1) = a_{11} e_1 + \ldots + a_{k1} e_k + 0 \cdot e_{k + 1} + \ldots + 0 \cdot e_n \\
            \vdots \\
            A(e_k) = a_{1k} e_1 + \ldots + a_{kk} e_k + 0 \cdot e_{k + 1} + \ldots + 0 \cdot e_n
        \end{cases}
    \end{equation*}
    Рассмотрим произвольный $x \in \mathcal{L}_1$. Выше уже было показано, что:
    \[x = x_1 e_1 + \ldots + x_k e_k + 0 \cdot e_{k + 1} + \ldots + 0 \cdot e_n\]
    Тогда из линейности преобразования $A$ получаем:
    \[A(x) = x_1 A(e_1) + \ldots + x_k A(e_k)\]
    Подставляя в равенство выше представления $A(e_1), \ldots, A(e_k)$, после приведения подобных слагаемых имеем некоторое разложение $A(x)$ по базису $\{e_1, \ldots, e_k\}$ подпространства $\mathcal{L}_1$, а значит $A(x) \in \mathcal{L}_1$. Но $x \in \mathcal{L}_1$ был выбран произвольно, откуда и получаем инвариантность $\mathcal{L}_1 = \left\langle e_1 , \ldots, e_k \right\rangle $ относительно линейного преобразования $A$, что и требовалось.
\end{proof}

\begin{theorem}
    Матрица $A$ линейного преобразования $A$ пространства $\mathcal{L}$ является \textbf{клеточно-диагональной} или \textbf{блочно-диагональной}, то есть имеет вид
    \begin{equation*}
        A = 
        \begin{pmatrix}
            A_1 & & & O \\
            & A_2 & & \\
            & & \ddots & \\
            O & & & A_s
        \end{pmatrix}
        ,
    \end{equation*}
    тогда и только тогда, когда базис $\mathcal{L}$ есть объединение базисов инвариантных относительно преобразования $A$ подпространств $\mathcal{L}_1, \ldots, \mathcal{L}_s$ (то есть $\mathcal{L} = \mathcal{L}_1 \oplus \ldots \oplus \mathcal{L}_s$, где $\mathcal{L}_i$ инвариантно относительно преобразования $A$).
\end{theorem}

\begin{proof}
    Требуемое легко следует из доказательства предыдущей теоремы посредством многократного применения содержащихся в нём рассуждений для каждого из подпространств, поэтому при желании смекалистый читатель без труда проведёт доказательство самостоятельно.
\end{proof}

\begin{lemma} \label{lemma:2.2}
    Если преобразования $A$ и $B$ перестановочны, то ядро и образ одного из них инвариантны относительно другого, то есть подпространства $\Ker A$ и $\Imm A$ инвариантны относительно $B$, а подпространства $\Ker B$ и $\Imm B$ инвариантны относительно $A$.
\end{lemma}

\begin{proof}
    Покажем сначала, что подпространство $\Ker A$ инвариантно относительно преобразования $B$. Возьмём произвольный $x \in \Ker A$. Тогда по определению ядра $Ax = 0$ (здесь $A$ понимается, естественным образом, как матрица преобразования $A$). Умножив это равенство на матрицу $B$ преобразования $B$ и вспомнив о перестановочности преобразований, имеем:
    \[BAx = 0 \iff ABx = 0 \iff A(Bx) = 0 \iff Bx \in \Ker A\]
    Итого получили, что $\forall x \in \Ker A \hookrightarrow B(x) \in \Ker A$. Это по определению и означает, что подпространство $\Ker A$ инвариантно относительно преобразования $B$.
    \\
    Увидим теперь, что и подпространство $\Imm A$ инвариантно относительно преобразования $B$. Вновь рассмотрим произвольный $x \in \Imm A$. Тогда по определению образа существует $z$ такой, что $A(z) = x$. Значит имеем следующее:
    \[Bx = BAz = [\text{в силу перестановочности}] = ABz = A(Bz) \in \Imm A\]
    Отсюда в силу произвольного выбора $x \in \Imm A$ получаем, что $\forall x \in \Imm A \hookrightarrow B(x) \in \Imm A$. Это вновь означает, что подпространство $\Imm A$ инвариантно относительно преобразования $B$.
    \\
    Осталось лишь увидеть, что доказательство нужного факта для подпространств $\Ker B$ и $\Imm B$ абсолютно аналогично, а значит мы можем с радостью заявить, что это победа!
\end{proof}

\subsection{Собственные подпространства и собственные векторы}

Мы найдём подпространство, инвариантное относительно заданного линейного преобразования $A$, если найдём преобразование, перестановочное с $A$ и имеющее ненулевое ядро (это прямое следствие леммы \ref{lemma:2.2}). Перестановочны с $A$ прежде всего многочлены от $A$ и, в частности, простейшие из них --- линейные. Умножив при необходимости линейный многочлен на число, напишем его в виде $A - \lambda E$, где $\lambda \in \mathbb{R}$.

\begin{definition}
    Если для числа $\lambda$ подпространство $\Ker(A - \lambda E)$ ненулевое, то $\lambda$ называется \textbf{собственным значением} преобразования, а подпространство --- \textbf{собственным подпространством}, отвечающим (или соответствующим) собственному значению $\lambda$.
\end{definition} \leavevmode
\\
Отметим один важный частный случай. Если преобразование $A$ имеет ненулевое ядро, то это ядро --- собственное подпространство, отвечающее собственному значению $\lambda = 0$. Ограничение $A$ на этом инвариантном подпространстве (то есть преобразование $A$ с суженной до $\Ker A$ областью определения) --- нулевое преобразование (тривиально следует из определения ядра). 
\\
Если ненулевой вектор $x$ лежит в собственном подпространстве, то для него выполнено:
\[(A - \lambda E)x = 0 \iff Ax = \lambda E x = \lambda x \iff Ax = \lambda x\]

\begin{lemma}
    Ограничение преобразования на собственном подпространстве является или нулевым преобразованием, или гомотетией: оно умножает каждый вектор этого подпространства на собственное значение. 
\end{lemma} \leavevmode
\\
Пусть нам каким-то образом удалось найти собственные значения преобразования $A$. Тогда для нахождения собственных подпространств нужно для каждого собственного значения $\lambda$ составить однородную систему линейных уравнений
\[(A - \lambda E)x = 0\]
Здесь $A$ --- матрица преобразования $A$ в некотором базисе $e$. Фундаментальная система решений этой ОСЛУ состоит из координатных столбцов векторов, составляющих базис собственного подпространства.

\begin{definition}
    \underline{Ненулевой} вектор $x$ назовём \textbf{собственным вектором} преобразования $A$, отвечающим собственному значению $\lambda$, если $Ax = \lambda x$.
\end{definition}

\begin{lemma}
    Собственные векторы и только они являются базисными векторами одномерных подпространств, инвариантных относительно $A$.
\end{lemma}

\begin{proof}
    $(\Longrightarrow):$ Пусть вектор $x$ собственный, а вектор $y$ принадлежит одномерному подпространству $\mathcal{L}_1$ с базисом $x$. Тогда $y = \alpha x$ и в силу линейности преобразования $A$:
    \[A(y) = \alpha A(x) = \alpha \lambda x = (\alpha \lambda) x \in \mathcal{L}_1\]
    Это и означает, что подпространство $\mathcal{L}_1$ инвариантно относительно преобразования $A$.
    \\
    $(\Longleftarrow):$ Пусть $x$ --- базис инвариантного подпространства $\mathcal{L}_1$. Тогда $A(x)$ лежит в $\mathcal{L}_1$ и раскладывается по базису: $A(x) = \lambda x$. Так как $x$ --- базис, то $x \neq 0$, а значит по определению он собственный.
\end{proof}

\begin{consequence}
    В собственном подпространстве через каждый вектор проходит одномерное инвариантное подпространство.
\end{consequence}

\begin{theorem}
    В $i$-ом столбце матрицы линейного преобразования все элементы вне главной диагонали равны нулю тогда и только тогда, когда $i$-ый базисный вектор собственный. В этом случае диагональный элемент столбца --- собственное значение.
\end{theorem}

\begin{proof}
    В самом деле, если базисный вектор $e_i$ собственный, то $A(e_i) = \lambda e_i$, и поэтому $i$-ый элемент координатного столбца вектора $A(e_i)$ равен $\lambda$, а остальные элементы равны нулю. Осталось лишь вспомнить, что координатный столбец $A(e_i)$ есть $i$-ый столбец матрицы преобразования. Обратное утверждение доказывается аналогично.
\end{proof}

\begin{remark}
    Если нам так чудесно повезло, что базис целиком состоит из собственных векторов, то матрица преобразования имеет диагональный вид с собственными значениями на диагонали. Этот факт есть легкое следствие теоремы выше.
\end{remark}

\subsection{Характеристическое уравнение}

Выберем базис и обозначим через $A$ матрицу линейного преобразования $A$ в этом базисе. Тогда преобразование $A - \lambda E$ имеет матрицу $A - \lambda E$, и по доказанной в предыдущей лекции теореме о совпадении размерности отображаемого пространства с суммой размерностей ядра отображения и образа отображения имеем, что $\Ker(A - \lambda E)$ отлично от нуля тогда и только тогда, когда 

\begin{equation} \label{eq:2.5}
    \det (A - \lambda E) = \det
    \begin{pmatrix}
        a_{11} - \lambda & a_{12} & \ldots & a_{1n} \\
        a_{21} & a_{22} - \lambda & \ldots & a_{2n} \\
        \vdots & \vdots & \ddots & \vdots \\
        a_{n1} & a_{n2} & \ldots & a_{nn} - \lambda
    \end{pmatrix}
    = 0
\end{equation} \leavevmode
\\
Равенство (\ref{eq:2.5}), рассматриваемое как условие на $\lambda$, называется \textbf{характеристическим уравнением} матрицы $A$, а его корни --- \textbf{характеристическими числами} матрицы $A$.

\begin{lemma}
    В вещественном пространстве все вещественные корни характеристического уравнения и только они являются собственными значениями.
\end{lemma} \leavevmode
\\
Левая часть характеристического уравнения представляет собой многочлен $p(\lambda)$ степени $n$. Действительно, этот факт легко следует из формулы полного разложения определителя. При этом очевидно, что $p(0) = \det A$ и что коэффициент при $\lambda^n$ равен $(-1)^n$, так как этот одночлен может получиться только из следующего единственного слагаемого единственным образом, а именно взятием из всех скобок $\lambda$:
\begin{equation*} \label{eq:staples:2.5} \tag{$\star$}
    (a_{11} - \lambda)(a_{22} - \lambda) \ldots (a_{nn} - \lambda)
\end{equation*}
Далее найдём коэффициент при $\lambda^{n - 1}$. Заметим, что в разложении определителя не может быть слагаемых, у которых ровно $(n - 1)$ сомножитель представляет собой скобку вида $(a_{ii} - \lambda)$, так как при выборе какого-либо набора из $(n - 1)$ диагонального элемента (они и есть скобки такого вида) оставшийся сомножитель определяется однозначно и является тем самым диагональным элементом, который ещё не выбрали (однозначность легко получается из того, что каждое слагаемое в формуле полного разложения определителя представляет собой такое произведение элементов матрицы, в котором заведомо есть по элементу из каждой строки и каждого столбца).
Это означает, что $\lambda^{n - 1}$ могло образоваться только при раскрытии скобок в том самом единственном слагаемом (\ref{eq:staples:2.5}). Тогда нетрудно видеть, поигравшись со скобочками и подобными слагаемыми, что коэффициент при $\lambda^{n - 1}$ будет в точности равен следу матрицы $A$, умноженному на $(-1)^{n - 1}$. Иначе говоря, искомый коэффициент есть следующее:
\[(-1)^{n - 1} (a_{11} + a_{22} + \ldots + a_{nn})\]
Таким образом,
\[p(\lambda) = (-1)^n \lambda^n + (-1)^{n - 1} (a_{11} + a_{22} + \ldots + a_{nn}) \lambda^{n - 1} + \ldots + \det A\]
Этот многочлен называется \textbf{характеристическим многочленом} матрицы $A$. Остальные его коэффициенты находить не будем, так как они нам не потребуются. 
\\
В вещественном пространстве может случится (при чётной размерности), что характеристическое уравнение не имеет ни одного вещественного корня, и, следовательно, линейное преобразование не имеет собственных значений и собственных подпространств. Примером может служить поворот плоскости на $30$ градусов.
\\
В вещественном пространстве нечётной размерности каждое линейное преобразование имеет хоть одно собственное значение и хоть одно собственное подпространство (это легко следует из того чудесного школьного факта, что многочлен нечётной степени всегда имеет хотя бы один вещественный корень).

\end{document}